\documentclass[a4paper,10pt]{article}
\usepackage{fontspec}
\usepackage{xeCJK}
\usepackage[a4paper,margin=20mm]{geometry}
\usepackage{graphicx}
\usepackage{fancyhdr}
\setsansfont[Path=/Users/yupyom/.src2tex/fonts/hackgen/, BoldFont=HackGen-Bold.ttf]{HackGen-Regular.ttf}
\setmonofont[Path=/Users/yupyom/.src2tex/fonts/hackgen/, BoldFont=HackGen-Bold.ttf]{HackGen-Regular.ttf}
\setCJKmainfont[Path=/usr/local/texlive/2026/texmf-dist/fonts/opentype/public/haranoaji/, Extension=.otf, BoldFont=HaranoAjiMincho-Bold.otf]{HaranoAjiMincho-Regular}
\setCJKsansfont[Path=/Users/yupyom/.src2tex/fonts/hackgen/, BoldFont=HackGen-Bold.ttf]{HackGen-Regular.ttf}
\setCJKmonofont[Path=/Users/yupyom/.src2tex/fonts/hackgen/, BoldFont=HackGen-Bold.ttf]{HackGen-Regular.ttf}
\XeTeXlinebreaklocale ""
\xeCJKsetup{CJKglue={},CJKecglue={}}
\newdimen\charwd
{\tt\global\setbox0=\hbox{x}\global\charwd=\wd0}
\frenchspacing
\pagestyle{fancy}
\renewcommand{\headrulewidth}{0pt}
\fancyhf{}
\fancyhead[R]{\rm src2tex version 2.236}
\fancyfoot[R]{\rm hanoi.m\qquad page \thepage}
% plain TeX compatibility
\makeatletter
\providecommand{\eqalign}[1]{\vcenter{\openup1\jot\m@th
  \ialign{\strut\hfil$\displaystyle{##}$&$\displaystyle{{}##}$\hfil\crcr#1\crcr}}}
\makeatother
\ifx\sevenrm\undefined
  \font\sevenrm=cmr7 scaled \magstep0
\fi
\def\mc{\relax}
\def\gt{\relax}
\def\sc{\scshape}
\begin{document}
\tt\mc 

\noindent
\mbox{}\rm\mc {\tt\%} \hrulefill
\tt\mc 

\noindent
\mbox{}\rm\mc {\tt\%} \ % beginning of TeX mode
\tt\mc 

\noindent
\mbox{}\rm\mc {\tt\%} \ \centerline{\bf Towers of Hanoi (Octave)}
\tt\mc 

\noindent
\mbox{}\rm\mc {\tt\%} 
\tt\mc 

\noindent
\mbox{}\rm\mc {\tt\%} \ \hspace{\leftmargini}This program gives an answer to the following famous problem (towers of
\tt\mc 

\noindent
\mbox{}\rm\mc {\tt\%} \ \hspace{\leftmargini}Hanoi).
\tt\mc 

\noindent
\mbox{}\rm\mc {\tt\%} \ \hspace{\leftmargini}There is a legend that when one of the temples in Hanoi was constructed,
\tt\mc 

\noindent
\mbox{}\rm\mc {\tt\%} \ \hspace{\leftmargini}three poles were erected and a tower consisting of 64 golden discs was
\tt\mc 

\noindent
\mbox{}\rm\mc {\tt\%} \ \hspace{\leftmargini}arranged on one pole, their sizes decreasing regularly from bottom to top.
\tt\mc 

\noindent
\mbox{}\rm\mc {\tt\%} \ \hspace{\leftmargini}The monks were to move the tower of discs to the opposite pole, moving
\tt\mc 

\noindent
\mbox{}\rm\mc {\tt\%} \ \hspace{\leftmargini}only one at a time, and never putting any size disc above a smaller one.
\tt\mc 

\noindent
\mbox{}\rm\mc {\tt\%} \ \hspace{\leftmargini}The job was to be done in the minimum numbers of moves. What strategy for
\tt\mc 

\noindent
\mbox{}\rm\mc {\tt\%} \ \hspace{\leftmargini}moving discs will accomplish this optimum transfer?
\tt\mc 

\noindent
\mbox{}\rm\mc {\tt\%} \ % end of TeX mode\hrulefill
\tt\mc 

\noindent
\mbox{}\rm\mc {\tt\%} \hrulefill\ hanoi.m \ \hrulefill
\tt\mc 

\noindent
\mbox{}\hfill

\noindent
\mbox{}global{\tt\mc \ }ARRAY;

\noindent
\mbox{}ARRAY{\tt\mc \ }={\tt\mc \ }8;{\tt\mc \ }{\tt\mc \ }{\tt\mc \ }{\tt\mc \ }{\tt\mc \ }{\tt\mc \ }{\tt\mc \ }{\tt\mc \ }{\tt\mc \ }{\tt\mc \ }{\tt\mc \ }{\tt\mc \ }{\tt\mc \ }{\tt\mc \ }{\tt\mc \ }{\tt\mc \ }{\tt\mc \ }{\tt\mc \ }{\tt\mc \ }{\tt\mc \ }{\tt\mc \ }{\tt\mc \ }{\tt\mc \ }{\tt\mc \ }{\tt\mc \ }{\tt\mc \ }{\tt\mc \ }{\tt\mc \ }{\tt\mc \ }{\tt\mc \ }\rm\mc {\tt\%} \ disc の数 \hfill
\tt\mc 

\noindent
\mbox{}\hfill

\noindent
\mbox{}global{\tt\mc \ }disc;

\noindent
\mbox{}disc{\tt\mc \ }={\tt\mc \ }zeros(3,{\tt\mc \ }ARRAY);{\tt\mc \ }{\tt\mc \ }{\tt\mc \ }{\tt\mc \ }{\tt\mc \ }{\tt\mc \ }{\tt\mc \ }{\tt\mc \ }{\tt\mc \ }{\tt\mc \ }{\tt\mc \ }{\tt\mc \ }{\tt\mc \ }{\tt\mc \ }{\tt\mc \ }{\tt\mc \ }{\tt\mc \ }\rm\mc {\tt\%} \ disc に関するデータの置き場所\hfill
\tt\mc 

\noindent
\mbox{}\hfill

\noindent
\mbox{}function{\tt\mc \ }init{\tt\_}array(){\tt\mc \ }{\tt\mc \ }{\tt\mc \ }{\tt\mc \ }{\tt\mc \ }{\tt\mc \ }{\tt\mc \ }{\tt\mc \ }{\tt\mc \ }{\tt\mc \ }{\tt\mc \ }{\tt\mc \ }{\tt\mc \ }{\tt\mc \ }{\tt\mc \ }{\tt\mc \ }{\tt\mc \ }{\tt\mc \ }{\tt\mc \ }\rm\mc {\tt\%} \ disc に関するデータの初期化\hfill
\tt\mc 

\noindent
\mbox{}{\tt\mc \ }{\tt\mc \ }{\tt\mc \ }{\tt\mc \ }global{\tt\mc \ }ARRAY{\tt\mc \ }disc;

\noindent
\mbox{}{\tt\mc \ }{\tt\mc \ }{\tt\mc \ }{\tt\mc \ }for{\tt\mc \ }j{\tt\mc \ }={\tt\mc \ }1:ARRAY

\noindent
\mbox{}{\tt\mc \ }{\tt\mc \ }{\tt\mc \ }{\tt\mc \ }{\tt\mc \ }{\tt\mc \ }{\tt\mc \ }{\tt\mc \ }disc(1,{\tt\mc \ }j){\tt\mc \ }={\tt\mc \ }ARRAY{\tt\mc \ }{\tt -}{\tt\mc \ }j{\tt\mc \ }+{\tt\mc \ }1;

\noindent
\mbox{}{\tt\mc \ }{\tt\mc \ }{\tt\mc \ }{\tt\mc \ }{\tt\mc \ }{\tt\mc \ }{\tt\mc \ }{\tt\mc \ }disc(2,{\tt\mc \ }j){\tt\mc \ }={\tt\mc \ }0;

\noindent
\mbox{}{\tt\mc \ }{\tt\mc \ }{\tt\mc \ }{\tt\mc \ }{\tt\mc \ }{\tt\mc \ }{\tt\mc \ }{\tt\mc \ }disc(3,{\tt\mc \ }j){\tt\mc \ }={\tt\mc \ }0;

\noindent
\mbox{}{\tt\mc \ }{\tt\mc \ }{\tt\mc \ }{\tt\mc \ }endfor

\noindent
\mbox{}endfunction

\noindent
\mbox{}\hfill

\noindent
\mbox{}global{\tt\mc \ }counter;

\noindent
\mbox{}counter{\tt\mc \ }={\tt\mc \ }0;{\tt\mc \ }{\tt\mc \ }{\tt\mc \ }{\tt\mc \ }{\tt\mc \ }{\tt\mc \ }{\tt\mc \ }{\tt\mc \ }{\tt\mc \ }{\tt\mc \ }{\tt\mc \ }{\tt\mc \ }{\tt\mc \ }{\tt\mc \ }{\tt\mc \ }{\tt\mc \ }{\tt\mc \ }{\tt\mc \ }{\tt\mc \ }{\tt\mc \ }{\tt\mc \ }{\tt\mc \ }{\tt\mc \ }{\tt\mc \ }{\tt\mc \ }{\tt\mc \ }{\tt\mc \ }{\tt\mc \ }\rm\mc {\tt\%} \ 移動回数カウンタ \hfill
\tt\mc 

\noindent
\mbox{}\hfill

\noindent
\mbox{}function{\tt\mc \ }print{\tt\_}result(){\tt\mc \ }{\tt\mc \ }{\tt\mc \ }{\tt\mc \ }{\tt\mc \ }{\tt\mc \ }{\tt\mc \ }{\tt\mc \ }{\tt\mc \ }{\tt\mc \ }{\tt\mc \ }{\tt\mc \ }{\tt\mc \ }{\tt\mc \ }{\tt\mc \ }{\tt\mc \ }{\tt\mc \ }\rm\mc {\tt\%} \ 結果の表示 \hfill
\tt\mc 

\noindent
\mbox{}{\tt\mc \ }{\tt\mc \ }{\tt\mc \ }{\tt\mc \ }global{\tt\mc \ }ARRAY{\tt\mc \ }disc{\tt\mc \ }counter;{\tt\mc \ }{\tt\mc \ }{\tt\mc \ }{\tt\mc \ }{\tt\mc \ }{\tt\mc \ }{\tt\mc \ }{\tt\mc \ }{\tt\mc \ }{\tt\mc \ }\rm\mc {\tt\%} \ Octave は 1-indexed のため\hfill
\tt\mc 

\noindent
\mbox{}{\tt\mc \ }{\tt\mc \ }{\tt\mc \ }{\tt\mc \ }counter{\tt\mc \ }={\tt\mc \ }counter{\tt\mc \ }+{\tt\mc \ }1;{\tt\mc \ }{\tt\mc \ }{\tt\mc \ }{\tt\mc \ }{\tt\mc \ }{\tt\mc \ }{\tt\mc \ }{\tt\mc \ }{\tt\mc \ }{\tt\mc \ }{\tt\mc \ }{\tt\mc \ }{\tt\mc \ }{\tt\mc \ }\rm\mc {\tt\%} \ ptr の値は実際の位置 $+1$ で管理する\hfill
\tt\mc 

\noindent
\mbox{}{\tt\mc \ }{\tt\mc \ }{\tt\mc \ }{\tt\mc \ }printf({\tt "}{\tt -}{\tt -}{\tt -}{\tt\#}{\tt\%}d{\tt -}{\tt -}{\tt -}{\tt\char92}n{\tt "},{\tt\mc \ }counter);

\noindent
\mbox{}{\tt\mc \ }{\tt\mc \ }{\tt\mc \ }{\tt\mc \ }for{\tt\mc \ }i{\tt\mc \ }={\tt\mc \ }1:3

\noindent
\mbox{}{\tt\mc \ }{\tt\mc \ }{\tt\mc \ }{\tt\mc \ }{\tt\mc \ }{\tt\mc \ }{\tt\mc \ }{\tt\mc \ }printf({\tt "}[{\tt\%}d]{\tt\mc \ }{\tt "},{\tt\mc \ }i{\tt\mc \ }{\tt -}{\tt\mc \ }1);

\noindent
\mbox{}{\tt\mc \ }{\tt\mc \ }{\tt\mc \ }{\tt\mc \ }{\tt\mc \ }{\tt\mc \ }{\tt\mc \ }{\tt\mc \ }for{\tt\mc \ }j{\tt\mc \ }={\tt\mc \ }1:ARRAY

\noindent
\mbox{}{\tt\mc \ }{\tt\mc \ }{\tt\mc \ }{\tt\mc \ }{\tt\mc \ }{\tt\mc \ }{\tt\mc \ }{\tt\mc \ }{\tt\mc \ }{\tt\mc \ }{\tt\mc \ }{\tt\mc \ }if{\tt\mc \ }disc(i,{\tt\mc \ }j){\tt\mc \ }!={\tt\mc \ }0

\noindent
\mbox{}{\tt\mc \ }{\tt\mc \ }{\tt\mc \ }{\tt\mc \ }{\tt\mc \ }{\tt\mc \ }{\tt\mc \ }{\tt\mc \ }{\tt\mc \ }{\tt\mc \ }{\tt\mc \ }{\tt\mc \ }{\tt\mc \ }{\tt\mc \ }{\tt\mc \ }{\tt\mc \ }printf({\tt "}{\tt\%}d{\tt\mc \ }{\tt "},{\tt\mc \ }disc(i,{\tt\mc \ }j));

\noindent
\mbox{}{\tt\mc \ }{\tt\mc \ }{\tt\mc \ }{\tt\mc \ }{\tt\mc \ }{\tt\mc \ }{\tt\mc \ }{\tt\mc \ }{\tt\mc \ }{\tt\mc \ }{\tt\mc \ }{\tt\mc \ }else

\noindent
\mbox{}{\tt\mc \ }{\tt\mc \ }{\tt\mc \ }{\tt\mc \ }{\tt\mc \ }{\tt\mc \ }{\tt\mc \ }{\tt\mc \ }{\tt\mc \ }{\tt\mc \ }{\tt\mc \ }{\tt\mc \ }{\tt\mc \ }{\tt\mc \ }{\tt\mc \ }{\tt\mc \ }break;

\noindent
\mbox{}{\tt\mc \ }{\tt\mc \ }{\tt\mc \ }{\tt\mc \ }{\tt\mc \ }{\tt\mc \ }{\tt\mc \ }{\tt\mc \ }{\tt\mc \ }{\tt\mc \ }{\tt\mc \ }{\tt\mc \ }endif

\noindent
\mbox{}{\tt\mc \ }{\tt\mc \ }{\tt\mc \ }{\tt\mc \ }{\tt\mc \ }{\tt\mc \ }{\tt\mc \ }{\tt\mc \ }endfor

\noindent
\mbox{}{\tt\mc \ }{\tt\mc \ }{\tt\mc \ }{\tt\mc \ }{\tt\mc \ }{\tt\mc \ }{\tt\mc \ }{\tt\mc \ }printf({\tt "}{\tt\char92}n{\tt "});

\noindent
\mbox{}{\tt\mc \ }{\tt\mc \ }{\tt\mc \ }{\tt\mc \ }endfor

\noindent
\mbox{}endfunction

\noindent
\mbox{}\hfill

\noindent
\mbox{}global{\tt\mc \ }ptr;

\noindent
\mbox{}ptr{\tt\mc \ }={\tt\mc \ }ones(1,{\tt\mc \ }3);{\tt\mc \ }{\tt\mc \ }{\tt\mc \ }{\tt\mc \ }{\tt\mc \ }{\tt\mc \ }{\tt\mc \ }{\tt\mc \ }{\tt\mc \ }{\tt\mc \ }{\tt\mc \ }{\tt\mc \ }{\tt\mc \ }{\tt\mc \ }{\tt\mc \ }{\tt\mc \ }{\tt\mc \ }{\tt\mc \ }{\tt\mc \ }{\tt\mc \ }{\tt\mc \ }{\tt\mc \ }{\tt\mc \ }\rm\mc {\tt\%} \ disc 移動用ポインタ（1-indexed）\hfill
\tt\mc 

\noindent
\mbox{}\hfill

\noindent
\mbox{}function{\tt\mc \ }move{\tt\_}one{\tt\_}disc(i,{\tt\mc \ }j){\tt\mc \ }{\tt\mc \ }{\tt\mc \ }{\tt\mc \ }{\tt\mc \ }{\tt\mc \ }{\tt\mc \ }{\tt\mc \ }{\tt\mc \ }{\tt\mc \ }{\tt\mc \ }{\tt\mc \ }\rm\mc {\tt\%} \ 1枚の disc を pole $i$ から\hfill
\tt\mc 

\noindent
\mbox{}{\tt\mc \ }{\tt\mc \ }{\tt\mc \ }{\tt\mc \ }{\tt\mc \ }{\tt\mc \ }{\tt\mc \ }{\tt\mc \ }{\tt\mc \ }{\tt\mc \ }{\tt\mc \ }{\tt\mc \ }{\tt\mc \ }{\tt\mc \ }{\tt\mc \ }{\tt\mc \ }{\tt\mc \ }{\tt\mc \ }{\tt\mc \ }{\tt\mc \ }{\tt\mc \ }{\tt\mc \ }{\tt\mc \ }{\tt\mc \ }{\tt\mc \ }{\tt\mc \ }{\tt\mc \ }{\tt\mc \ }{\tt\mc \ }{\tt\mc \ }{\tt\mc \ }{\tt\mc \ }{\tt\mc \ }{\tt\mc \ }{\tt\mc \ }{\tt\mc \ }{\tt\mc \ }{\tt\mc \ }{\tt\mc \ }{\tt\mc \ }\rm\mc {\tt\%} \ pole $j$ に移動する \hfill
\tt\mc 

\noindent
\mbox{}{\tt\mc \ }{\tt\mc \ }{\tt\mc \ }{\tt\mc \ }global{\tt\mc \ }disc{\tt\mc \ }ptr;

\noindent
\mbox{}{\tt\mc \ }{\tt\mc \ }{\tt\mc \ }{\tt\mc \ }ptr(i){\tt\mc \ }={\tt\mc \ }ptr(i){\tt\mc \ }{\tt -}{\tt\mc \ }1;

\noindent
\mbox{}{\tt\mc \ }{\tt\mc \ }{\tt\mc \ }{\tt\mc \ }disc(j,{\tt\mc \ }ptr(j)){\tt\mc \ }={\tt\mc \ }disc(i,{\tt\mc \ }ptr(i));

\noindent
\mbox{}{\tt\mc \ }{\tt\mc \ }{\tt\mc \ }{\tt\mc \ }ptr(j){\tt\mc \ }={\tt\mc \ }ptr(j){\tt\mc \ }+{\tt\mc \ }1;

\noindent
\mbox{}{\tt\mc \ }{\tt\mc \ }{\tt\mc \ }{\tt\mc \ }disc(i,{\tt\mc \ }ptr(i)){\tt\mc \ }={\tt\mc \ }0;

\noindent
\mbox{}endfunction

\noindent
\mbox{}\hfill

\noindent
\mbox{}function{\tt\mc \ }move{\tt\_}discs(n,{\tt\mc \ }i,{\tt\mc \ }j,{\tt\mc \ }k){\tt\mc \ }{\tt\mc \ }{\tt\mc \ }{\tt\mc \ }{\tt\mc \ }{\tt\mc \ }{\tt\mc \ }{\tt\mc \ }{\tt\mc \ }\rm\mc {\tt\%} \ \underline{\textsf{上から $n$ 枚目までの disc}}を、\hfill
\tt\mc 

\noindent
\mbox{}{\tt\mc \ }{\tt\mc \ }{\tt\mc \ }{\tt\mc \ }{\tt\mc \ }{\tt\mc \ }{\tt\mc \ }{\tt\mc \ }{\tt\mc \ }{\tt\mc \ }{\tt\mc \ }{\tt\mc \ }{\tt\mc \ }{\tt\mc \ }{\tt\mc \ }{\tt\mc \ }{\tt\mc \ }{\tt\mc \ }{\tt\mc \ }{\tt\mc \ }{\tt\mc \ }{\tt\mc \ }{\tt\mc \ }{\tt\mc \ }{\tt\mc \ }{\tt\mc \ }{\tt\mc \ }{\tt\mc \ }{\tt\mc \ }{\tt\mc \ }{\tt\mc \ }{\tt\mc \ }{\tt\mc \ }{\tt\mc \ }{\tt\mc \ }{\tt\mc \ }{\tt\mc \ }{\tt\mc \ }{\tt\mc \ }{\tt\mc \ }\rm\mc {\tt\%} \ pole $i$ から pole $j$ に\hfill
\tt\mc 

\noindent
\mbox{}{\tt\mc \ }{\tt\mc \ }{\tt\mc \ }{\tt\mc \ }{\tt\mc \ }{\tt\mc \ }{\tt\mc \ }{\tt\mc \ }{\tt\mc \ }{\tt\mc \ }{\tt\mc \ }{\tt\mc \ }{\tt\mc \ }{\tt\mc \ }{\tt\mc \ }{\tt\mc \ }{\tt\mc \ }{\tt\mc \ }{\tt\mc \ }{\tt\mc \ }{\tt\mc \ }{\tt\mc \ }{\tt\mc \ }{\tt\mc \ }{\tt\mc \ }{\tt\mc \ }{\tt\mc \ }{\tt\mc \ }{\tt\mc \ }{\tt\mc \ }{\tt\mc \ }{\tt\mc \ }{\tt\mc \ }{\tt\mc \ }{\tt\mc \ }{\tt\mc \ }{\tt\mc \ }{\tt\mc \ }{\tt\mc \ }{\tt\mc \ }\rm\mc {\tt\%} \ pole $k$ を経由して、移動する\hfill
\tt\mc 

\noindent
\mbox{}{\tt\mc \ }{\tt\mc \ }{\tt\mc \ }{\tt\mc \ }if{\tt\mc \ }n{\tt\mc \ }{\tt >}={\tt\mc \ }1

\noindent
\mbox{}{\tt\mc \ }{\tt\mc \ }{\tt\mc \ }{\tt\mc \ }{\tt\mc \ }{\tt\mc \ }{\tt\mc \ }{\tt\mc \ }move{\tt\_}discs(n{\tt\mc \ }{\tt -}{\tt\mc \ }1,{\tt\mc \ }i,{\tt\mc \ }k,{\tt\mc \ }j);{\tt\mc \ }{\tt\mc \ }{\tt\mc \ }{\tt\mc \ }{\tt\mc \ }\rm\mc {\tt\%} \ 関数 {\tt move\_discs()}の中で、さらに自分自身 \hfill
\tt\mc 

\noindent
\mbox{}{\tt\mc \ }{\tt\mc \ }{\tt\mc \ }{\tt\mc \ }{\tt\mc \ }{\tt\mc \ }{\tt\mc \ }{\tt\mc \ }move{\tt\_}one{\tt\_}disc(i,{\tt\mc \ }j);{\tt\mc \ }{\tt\mc \ }{\tt\mc \ }{\tt\mc \ }{\tt\mc \ }{\tt\mc \ }{\tt\mc \ }{\tt\mc \ }{\tt\mc \ }{\tt\mc \ }{\tt\mc \ }{\tt\mc \ }\rm\mc {\tt\%} \ {\tt move\_discs()} が使われている。このような \hfill
\tt\mc 

\noindent
\mbox{}{\tt\mc \ }{\tt\mc \ }{\tt\mc \ }{\tt\mc \ }{\tt\mc \ }{\tt\mc \ }{\tt\mc \ }{\tt\mc \ }print{\tt\_}result();{\tt\mc \ }{\tt\mc \ }{\tt\mc \ }{\tt\mc \ }{\tt\mc \ }{\tt\mc \ }{\tt\mc \ }{\tt\mc \ }{\tt\mc \ }{\tt\mc \ }{\tt\mc \ }{\tt\mc \ }{\tt\mc \ }{\tt\mc \ }{\tt\mc \ }{\tt\mc \ }{\tt\mc \ }\rm\mc {\tt\%} \ 手法は、「再帰的呼びだし」といわれる。 \hfill
\tt\mc 

\noindent
\mbox{}{\tt\mc \ }{\tt\mc \ }{\tt\mc \ }{\tt\mc \ }{\tt\mc \ }{\tt\mc \ }{\tt\mc \ }{\tt\mc \ }move{\tt\_}discs(n{\tt\mc \ }{\tt -}{\tt\mc \ }1,{\tt\mc \ }k,{\tt\mc \ }j,{\tt\mc \ }i);

\noindent
\mbox{}{\tt\mc \ }{\tt\mc \ }{\tt\mc \ }{\tt\mc \ }endif

\noindent
\mbox{}endfunction

\noindent
\mbox{}\hfill

\noindent
\mbox{}\rm\mc {\tt\%} \par\begin{center}\includegraphics[scale=0.3]{hanoi1}\quad\includegraphics[scale=0.3]{hanoi2}\end{center}
\tt\mc 

\noindent
\mbox{}\rm\mc {\tt\%}
\tt\mc 

\noindent
\mbox{}\rm\mc {\tt\%} \ たとえば、関数 {\tt move\_discs(4, 1, 2, 3)} を呼び出すことは、
\tt\mc 

\noindent
\mbox{}\rm\mc {\tt\%} \ 上図のような操作をすることに対応する。\hfill
\tt\mc 

\noindent
\mbox{}\hfill

\noindent
\mbox{}ptr(1){\tt\mc \ }={\tt\mc \ }ARRAY{\tt\mc \ }+{\tt\mc \ }1;{\tt\mc \ }{\tt\mc \ }{\tt\mc \ }{\tt\mc \ }{\tt\mc \ }{\tt\mc \ }{\tt\mc \ }{\tt\mc \ }{\tt\mc \ }{\tt\mc \ }{\tt\mc \ }{\tt\mc \ }{\tt\mc \ }{\tt\mc \ }{\tt\mc \ }{\tt\mc \ }{\tt\mc \ }{\tt\mc \ }{\tt\mc \ }{\tt\mc \ }{\tt\mc \ }\rm\mc {\tt\%} \ 1-indexed のため ARRAY $+1$ \hfill
\tt\mc 

\noindent
\mbox{}ptr(2){\tt\mc \ }={\tt\mc \ }1;

\noindent
\mbox{}ptr(3){\tt\mc \ }={\tt\mc \ }1;

\noindent
\mbox{}\hfill

\noindent
\mbox{}init{\tt\_}array();

\noindent
\mbox{}move{\tt\_}discs(ARRAY,{\tt\mc \ }1,{\tt\mc \ }2,{\tt\mc \ }3);{\tt\mc \ }{\tt\mc \ }{\tt\mc \ }{\tt\mc \ }{\tt\mc \ }{\tt\mc \ }{\tt\mc \ }{\tt\mc \ }{\tt\mc \ }{\tt\mc \ }{\tt\mc \ }{\tt\mc \ }{\tt\mc \ }\rm\mc {\tt\%} \ {\tt ARRAY} 枚の disc をpole 0 から pole 1 に pole 2\hfill
\tt\mc 

\noindent
\mbox{}{\tt\mc \ }{\tt\mc \ }{\tt\mc \ }{\tt\mc \ }{\tt\mc \ }{\tt\mc \ }{\tt\mc \ }{\tt\mc \ }{\tt\mc \ }{\tt\mc \ }{\tt\mc \ }{\tt\mc \ }{\tt\mc \ }{\tt\mc \ }{\tt\mc \ }{\tt\mc \ }{\tt\mc \ }{\tt\mc \ }{\tt\mc \ }{\tt\mc \ }{\tt\mc \ }{\tt\mc \ }{\tt\mc \ }{\tt\mc \ }{\tt\mc \ }{\tt\mc \ }{\tt\mc \ }{\tt\mc \ }{\tt\mc \ }{\tt\mc \ }{\tt\mc \ }{\tt\mc \ }{\tt\mc \ }{\tt\mc \ }{\tt\mc \ }{\tt\mc \ }{\tt\mc \ }{\tt\mc \ }{\tt\mc \ }{\tt\mc \ }\rm\mc {\tt\%} \ を経由して、移動する \hfill

\tt\mc 

\rm\mc

\end{document}
